\documentclass[journal]{IEEEtran}
\usepackage{amsmath,amssymb,mathtools}
\usepackage{amsthm}
\usepackage{booktabs,array}
\usepackage{enumitem}
\usepackage{microtype}
\usepackage[hidelinks]{hyperref}
\newtheoremstyle{plainn}{}{}{\itshape}{}{\bfseries}{.}{.5em}%
  {\thmname{#1}\thmnumber{ #2}\thmnote{ (#3)}}
\theoremstyle{plainn}
\newtheorem{thm}{Theorem}
\newtheorem{prop}[thm]{Proposition}
\newtheorem{lem}[thm]{Lemma}
\newtheorem{cor}[thm]{Corollary}
\newtheorem{claim}{Claim}
\newtheorem{defn}[thm]{Definition}
\theoremstyle{remark}
\newtheorem{remark}[thm]{Remark}
\newtheorem{example}[thm]{Example}
\renewcommand{\qedsymbol}{QED}

\setlength{\abovedisplayskip}{5pt plus 2pt minus 2pt}
\setlength{\belowdisplayskip}{5pt plus 2pt minus 2pt}
\setlength{\abovedisplayshortskip}{3pt plus 1pt}
\setlength{\belowdisplayshortskip}{4pt plus 2pt minus 2pt}

\begin{document}
\title{The free-row fibre: a polynomial criterion and the geometry of
functional observer feasibility}
\author{}
\maketitle

\begin{abstract}
For a linear system with measured output and a prescribed block $L_0$ of observer rows, a functional
observer of order $q$ exists exactly when some matrix $L=[L_0;R]$ satisfies an algebraic and a
spectral rank condition. Whether such an $R$ exists is a quantifier alternation: one asks for a
subspace of observer row space for which a certain family of matrices has constant rank. This paper
studies that family directly. Fixing the compressed row space $U$ of an observer and passing to the
free rows, the admissible values of $LA$ form an affine space --- the \emph{fibre} --- whose dimension
is $q\,a$, where $a=\dim(U\cap T)$ and $T$ is the intrinsic space of row pairs attainable by rows
that vanish on the unobservable subspace. On the fibre the spectral condition is equivalent to the
vanishing of the degree of a single greatest common divisor, computed over the rational function field
of the fibre; the criterion is polynomial for every fixed compressed rank and does not enumerate
maximal minors. We then characterise when the whole fibre fails. If the number $q$ of observer rows
exceeds the compressed rank $d'$ and the fibre direction is common to all rows --- which is the case
here --- then the failure of the fibre is necessarily a fixed spectral factor: the common divisor is
independent of the fibre parameters. Consequently the bad set is closed in the Grassmannian of
admissible $U$, the feasible set is relatively open, and infeasibility is exactly the statement that
the whole compatibility stratum consists of bad points. The failure test reduces to checking whether a
certain resultant vanishes identically, which gives a one-sided (randomised) decision procedure; on an
explicit parameter region --- when compatibility is automatic --- the outer existence question
inherits it. We also record what the criterion does \emph{not} settle: the thin stratum where
compatibility is a genuine constraint remains open, and no hardness result is claimed.
\end{abstract}

\begin{IEEEkeywords}
functional observer; polynomial matrix; greatest common divisor; Grassmannian; constant rank;
randomised decision procedure.
\end{IEEEkeywords}

\section{Introduction}\label{sec:intro}

A functional observer estimates a prescribed linear function of the state without estimating the
state itself, and its order can be far below the order of a Luenberger observer. The design question
is quantitative: for a given system and a given set of functionals, which orders are attainable, and
can that be decided cheaply? The characterization is classical \cite{fer10,dar25} and has recently
been completed order by order \cite{fer26a,fer26b}; the computational question is not settled. The
companion paper \cite{paperI} gives parameterized upper bounds for deciding whether an observer of
prescribed order exists; the companion paper \cite{paperII} determines the minimum order. Both work
with the compressed object introduced below, but neither analyses the object this paper is about.

The two observer conditions couple the observer matrix and its product with the system matrix: the
same free rows produce both. Write $\mathcal L=[L_0;R]$ and $L$ for $\mathcal L$. With
$K_0=\ker[C;CA]$ and a basis $N$ of $K_0$, the conditions are

\begin{equation}\label{eq:fivea}
\operatorname{rank}[\mathcal LA;\mathcal L;CA;C]=\operatorname{rank}[\mathcal L;CA;C],
\end{equation}
\begin{equation}\label{eq:fiveb}
\operatorname{rank}[\lambda\mathcal L-\mathcal LA;CA;C]=\operatorname{rank}[\mathcal L;CA;C]
\quad\text{for all }\lambda .
\end{equation}
Both depend on $L$ only through the two compressed blocks
\begin{equation}\label{eq:HHA}
H=LN\in\mathbb C^{q\times k},\qquad H_A=LAN\in\mathbb C^{q\times k},
\end{equation}
and become, respectively, a row-space inclusion and a constant-rank condition:
\begin{equation}\label{eq:compressed}
\operatorname{row}H_A\subseteq\operatorname{row}H,\qquad
\operatorname{rank}(\lambda H-H_A)=\operatorname{rank}H\quad\forall\lambda .
\end{equation}
The pair $(H,H_A)$ is not arbitrary: it ranges over the row pairs attainable by $r$ free rows, a
subspace of $\mathbb C^{2k}$ that depends only on $A$, $C$ and $L_0$. Choosing those rows is the
only design freedom, and the two conditions pull in opposite directions --- $(5a)$ wants $H_A$
small, $(5b)$ wants the pencil $\lambda H-H_A$ to keep its rank.

This paper fixes $H$ and asks what happens to $(5b)$ as $H_A$ runs over the values the free rows can
give it. The answer is an affine space, the fibre, and on that space the spectral condition is a
single polynomial-degree condition. The three facts that make this useful are that the fibre is an
affine space of known dimension, that its dimension is governed by $a=\dim(U\cap T)$ --- with $U$ the
row space of $H$ and $T$ the intrinsic space of row pairs --- and that the failure of $(5b)$ on the
whole fibre cannot be parameter-dependent when $q>d'$, the number of rows exceeding the compressed
rank.

\subsection{What is proved}

Section~\ref{sec:prelim} sets up the fibre and proves that it is empty or affine of dimension $q\,a$
(Theorem~\ref{thm:fibre}). Section~\ref{sec:criterion} proves the criterion: on the fibre,
$(5b)$ holds for some point if and only if the greatest common divisor of all $d'$-minors of
$\lambda H-H_A$, computed over the rational function field of the fibre, has degree zero
(Theorem~\ref{thm:criterion}); the criterion depends only on the intrinsic attainable set
(Theorem~\ref{thm:intrinsic}), the witness it produces has polynomial height
(Theorem~\ref{thm:witness}), and testing it is a one-sided randomised procedure
(Theorem~\ref{thm:rp}). Section~\ref{sec:bad} characterises the failure of the whole fibre: when
$q>d'$ the failure is a fixed spectral factor (Theorem~\ref{thm:factor}), hence the bad set is closed
(Theorem~\ref{thm:closed}) and infeasibility is exactly containment of the compatibility stratum in
the bad set (Theorem~\ref{thm:no}). Section~\ref{sec:algorithms} turns this into procedures on an
explicit region (Theorem~\ref{thm:regime}) and settles the stratum $t=1$ by enumeration
(Theorem~\ref{thm:t1}). Section~\ref{sec:open} states what remains open.

\subsection{What is not claimed}

No hardness result is claimed, and no polynomial-time decision procedure for the general case. The
criterion is polynomial for fixed compressed rank, and the randomised procedure decides the outer
question only where the compatibility stratum is large enough to be sampled; the thin stratum where
compatibility is a genuine constraint is left open (Section~\ref{sec:open}). We also do not claim that
the fibre criterion is \emph{the} natural algorithm: it is the criterion that the geometry of the
fibre supplies.

\section{Preliminaries}\label{sec:prelim}

\subsection{Notation}\label{sec:notation}

Matrices have complex entries and rational data; $\operatorname{row}M$ is the row space and
$\operatorname{rank}M$ its dimension. For a subspace $U$ and a matrix $X$ with $k$ columns,
$UX=\{uX:u\in U\}$ is a subspace, and $U-X=\{u-X_i:\text{row }i\}$ is read row by row below. A pencil
is a matrix whose entries are polynomials of degree at most one in $\lambda$. The Grassmannian of
$d'$-dimensional subspaces of $\mathbb C^k$ is $\mathrm{Gr}(d',k)$; a subspace containing a fixed
$U_0$ is written $U\supseteq U_0$. For a field $F$ and a set $S\subseteq F[\lambda]$, $\gcd S$ is the
monic greatest common divisor, with $\gcd\varnothing=1$ by convention.

\subsection{Attainable row pairs and the fibre}\label{sec:attainable}

\begin{lem}\label{lem:pair}
The pair $(H_i,(H_A)_i)$ of row $i$ ranges over $\operatorname{row}[N\ |\ AN]$ as $L$ ranges over the
$q\times n$ matrices with $LN=H$ and $i$ over the rows. Consequently, with $U=\operatorname{row}H$
and
\begin{equation}\label{eq:Kset}
\mathcal K(H)=\bigl\{K\in\mathbb C^{q\times k}:\ (H_i,K_i)\in\operatorname{row}[N\ |\ AN]\ \
\forall i\bigr\},
\end{equation}
$H_A$ ranges over $\mathcal K(H)$, and $(5a)$ is the condition $K\in\mathcal K(H)\cap U^{q}$, i.e.\
$K_i\in U$ for every $i$.
\end{lem}

\begin{proof}
A row $l\in\mathbb C^{n}$ yields the pair $(lN,lAN)\in\operatorname{row}[N\ |\ AN]$; conversely every
element of that row space is of this form because its first half determines $l$ on $K_0$ and $N$ is a
basis of $K_0$. Gather the rows. The last claim is the definition of $(5a)$ in
(\ref{eq:compressed}).
\end{proof}

To describe the fibre we use the space of row pairs that vanish on $K_0$. Put
\begin{equation}\label{eq:T}
T=\{lAN:\ lN=0\}\subseteq\mathbb C^{k}.
\end{equation}
This space depends only on $A$ and $K_0$, not on any choice of coordinates: if $lN=0$ then
$lNZ'=0$ for every $Z'$, and $T=\{l(AN-NZ'):lN=0\}$. Choose $Z'\in\mathbb C^{k\times k}$ and write
$W=AN-NZ'$; then $(lN,lAN)=(lN,\,lNZ'+lW)$, so
\begin{equation}\label{eq:decomp}
H_A=HZ''+G',\qquad Z''=Z'+W_{\mathrm{top}},\qquad G'\in T^{q},
\end{equation}
where $W_{\mathrm{top}}$ collects the rows of $W$ that are not in $T$'s complement position; the
decomposition is the standard one of the compressed pencil. Different choices of $Z'$ translate
$G'$ inside $T^{q}$ and leave the pencil $\lambda H-H_A$ unchanged.

\begin{thm}[Fibre]\label{thm:fibre}
Fix $H$ with $U=\operatorname{row}H$ and $d'=\operatorname{rank}H$. The set of admissible $H_A$ for
this $H$ is $HZ''+\mathcal F_U$ with
\begin{equation}\label{eq:fibre}
\mathcal F_U=\prod_{i=1}^{q}\bigl(T\cap(U-H_iZ'')\bigr).
\end{equation}
Either $\mathcal F_U=\varnothing$, or $\mathcal F_U$ is an affine space of dimension $q\,a$ with
$a=\dim(U\cap T)$. Moreover $\mathcal F_U\ne\varnothing$ if and only if
$UZ''\subseteq U+T$, and then $\mathcal F_U\cong\prod_i (U\cap T)$.
\end{thm}

\begin{proof}
By Lemma~\ref{lem:pair}, $H_A=HZ''+G'$ with $G'\in T^{q}$ is admissible for $H$ exactly when
$H_iZ''+G_i'\in U$ for every $i$, which is $G_i'\in T\cap(U-H_iZ'')$. The set $T\cap(U-H_iZ'')$ is
an intersection of an affine space with a subspace, hence empty or a coset of $T\cap U$. Taking
products gives the dimension $q\cdot\dim(T\cap U)$. For nonemptiness, $T\cap(U-x)\ne\varnothing$ iff
$x\in U+T$, and rows of $H_iZ''$ with $H_i\in U$ satisfy this for all $i$ exactly when
$UZ''\subseteq U+T$.
\end{proof}

\begin{remark}\label{rem:T-invariant}
$T$ is intrinsic but $Z''$ and $T$'s description change with the choice of complement in
$\mathbb C^n$; Theorem~\ref{thm:intrinsic} below is the statement that the criterion does not depend
on that choice. Everything is computed from $(H,Z'',T)$, and different decompositions give the same
fibre up to translation inside $T^{q}$.
\end{remark}

\section{The criterion}\label{sec:criterion}

Throughout, $F=\mathbb C(t)$ is the rational function field of the fibre parameters, and for
$G'\in F^{q\times k}$ we write $P(\lambda)=\lambda H-HZ''-G'\in F[\lambda]^{q\times k}$. Only
$(\lambda H-H_A)$ matters, so the criterion below is stated for a general $K=HZ''+G'\in F^{q\times
k}$ with $G'\in T^{q}$.

\begin{thm}[Criterion]\label{thm:criterion}
Assume $\mathcal F_U\ne\varnothing$. The following are equivalent.
\begin{enumerate}[label=(\roman*),leftmargin=2.2em,itemsep=1pt]
\item There exists $G'\in\mathcal F_U$ such that $\operatorname{rank}P(\lambda)=d'$ for all
$\lambda$.
\item $\deg_\lambda g=0$, where $g$ is the greatest common divisor in $F[\lambda]$ of all $d'\times
d'$ minors of $\lambda H-HZ''-G'$.
\end{enumerate}
\end{thm}

\begin{proof}
Write $f_j$ for the minors. A matrix over $F[\lambda]$ has constant rank $d'$ for every
$\lambda\in\overline F$ if and only if the ideal generated by the $f_j$ contains a nonzero constant,
i.e.\ $\gcd_j f_j\in F^{\times}$; this is the standard gcd criterion for a polynomial matrix to have
no eigenvalue \cite[Ch.~6]{vzgg13}. For the equivalence with the existence of a point, note that the
$f_j$ are polynomials in the fibre parameters $t$ with coefficients in $\mathbb C[\lambda]$, so
$\deg g=0$ says that some finite $\mathbb Z$-combination of the $f_j$ (after clearing denominators)
is a nonzero constant $D\in\mathbb C[t]$. Evaluating at $t_0$ with $D(t_0)\ne0$ gives a point of the
affine space $\mathcal F_U$ at which the specialised minors have gcd one, hence constant rank $d'$;
the converse is immediate because a common root at one point would make $\gcd$ positive at the
generic point. Since $\mathcal F_U$ is affine over $\mathbb C$ with the coordinates $t$, and $D$ is a
nonzero polynomial, such a $t_0$ exists.
\end{proof}

The step that makes the criterion effective is that a positive-degree common divisor cannot
disappear under specialisation, provided its leading coefficient in $\lambda$ is a constant. That is
the only place where the minors' structure is used, and it is machine-checked.

\begin{lem}[Specialisation]\label{lem:specialize}
Let $h\in\mathbb C[t][\lambda]$ have $\lambda$-leading coefficient equal to a nonzero constant $c$.
Then for every $t_0$, $\operatorname{natDegree}\bigl(h(t_0,\cdot)\bigr)=\operatorname{natDegree}h$;
consequently $h\mid f$ in $\mathbb C[t][\lambda]$ implies $h(t_0,\cdot)\mid f(t_0,\cdot)$.
\end{lem}

\begin{proof}
Write $h=\sum_{j\le n}h_j\lambda^{j}$ with $h_n=c$. The coefficient of $\lambda^{n}$ in
$h(t_0,\cdot)$ is $h_n(t_0)=c\ne0$, so the degree cannot drop; for the divisibility, apply the ring
homomorphism $\mathbb C[t]\to\mathbb C$, $t\mapsto t_0$, to the identity $f=qh$.
\end{proof}

The leading coefficient of at least one $d'$-minor is a constant of this kind: the minor on rows $I$
and columns $J$ with $\det H_{IJ}\ne0$ has leading coefficient $\det H_{IJ}$, nonzero. Hence a
positive-degree $g$ has a primitive representative with constant leading coefficient, and
Lemma~\ref{lem:specialize} applies: every point of the fibre is bad. This is the content of
Corollary~\ref{cor:every-bad} below.

\begin{thm}[Intrinsic form]\label{thm:intrinsic}
The criterion of Theorem~\ref{thm:criterion} depends only on $H$ and on the attainable set
$\mathcal K(H)$ of \eqref{eq:Kset}, not on the choice of $Z'$, of $T$'s complement, or of the basis
of $K_0$: feasibility is the existence of $K\in\mathcal K(H)\cap U^{q}$ such that the $d'$-minors of
$\lambda H-K$ have gcd one.
\end{thm}

\begin{proof}
Both objects $\mathcal K(H)$ and $T$ are defined from $A$, $C$ and $K_0$ alone
(Lemma~\ref{lem:pair}, \eqref{eq:T}), and the map $K\mapsto G'=K-HZ''$ is a bijection between
$\mathcal K(H)\cap U^{q}$ and $\mathcal F_U$ by Theorem~\ref{thm:fibre}. A bijection between
parametrisations does not change the function field, so the gcd is computed in the same field and its
degree is the same.
\end{proof}

\begin{thm}[Witness]\label{thm:witness}
Suppose $U$ is defined over $\mathbb Q$, $\mathcal F_U\ne\varnothing$ and $\deg_\lambda g=0$. Then
there is $G'\in\mathcal F_U(\mathbb Q)$ whose entries have bit size
$\operatorname{poly}(\mathcal L,d')$, where $\mathcal L$ is the input length.
\end{thm}

\begin{proof}
Take two constant-coefficient combinations $F_1,F_2$ of the minors with
$\gcd_{F[\lambda]}(F_1,F_2)=1$ (generic combinations work, since the total gcd is one). Their
resultant $D=\operatorname{Res}_\lambda(F_1,F_2)\in\mathbb Q[t]$ is nonzero, and the adjugate of the
Sylvester matrix gives $UF_1+VF_2=D$ with polynomial $U,V$, so no clearing of denominators is
needed. The Sylvester matrix has order at most $2d'$ and entries of $t$-degree at most $d'$, hence
$\deg_t D\le 2d'^{2}$. A nonzero polynomial of degree $\delta$ has a point of the grid
$\{0,\dots,\delta\}^{N}$ where it does not vanish (induction on $N$), so a rational $t_0$ with
$D(t_0)\ne0$ exists with coordinates of size $O(\log\delta)$. At such a point the specialised minors
have gcd one: any common divisor $h$ would divide $F_1(t_0),F_2(t_0)$ and hence $D(t_0)\ne0$. The
fibre coordinates are affine in $t_0$, and the free rows are recovered from $K$ by solving
$l[N\ AN]=K_i$, which adds only polynomial height.
\end{proof}

\begin{thm}[One-sided test]\label{thm:rp}
Let $F_1,F_2$ be generic constant-coefficient combinations of the minors and
$D=\operatorname{Res}_\lambda(F_1,F_2)$. Then the fibre is entirely bad if and only if $D\equiv0$.
Consequently, evaluating $D$ at a random rational point $t_0$ and declaring the fibre feasible when
the value is nonzero is a one-sided procedure: it never reports a bad fibre as feasible, and reports
a feasible fibre with probability at least $1-2d'^{2}/|S|$ for a grid $S$.
\end{thm}

\begin{proof}
If every point of the fibre is bad then for every $t$ the specialised minors have a common root, so
$F_1(t,\cdot)$ and $F_2(t,\cdot)$ have one, hence $D(t)=0$. If the fibre is feasible, then
$\deg g=0$ and for generic combinations $\gcd(F_1,F_2)=1$, so $D\not\equiv0$; the bound is
Schwartz--Zippel applied to $D$ of degree at most $2d'^{2}$.
\end{proof}

\section{The bad set}\label{sec:bad}

\begin{defn}\label{def:bad}
$U\in\mathrm{Gr}(d',k)$ is \emph{bad} if $\mathcal F_U\ne\varnothing$ and the fibre is entirely bad,
i.e.\ every $G'\in\mathcal F_U$ fails $(5b)$. Write $\mathcal B$ for the bad set.
\end{defn}

\begin{cor}\label{cor:every-bad}
If $\deg_\lambda g\ge1$ then every point of the fibre is bad: $g$'s primitive representative has
constant leading coefficient and divides every minor, so Lemma~\ref{lem:specialize} preserves a
common factor at every $t_0$.
\end{cor}

The next theorem is the structural statement this paper is built around: the failure cannot be
parameter-dependent once $q>d'$, because the fibre direction is common to all rows.

\begin{thm}[Fixed spectral factor]\label{thm:factor}
Let $U$ be such that $\mathcal F_U\ne\varnothing$, $q>d'$, and $M$ be a matrix with
$H_A=HZ''+Y'M$, $Y'\in\mathbb C^{q\times a}$, $M\in\mathbb C^{a\times d'}$ (the fibre direction is
common to all rows, as in \eqref{eq:fibre}). If the fibre is entirely bad then $g\in\mathbb C[\lambda]$:
the common divisor does not depend on the fibre parameters.
\end{thm}

\begin{proof}
Write $P(Y,\lambda)=\lambda X-B_0-YM$ with $X=H$, $B_0=HZ''$. Consider the incidence variety
\begin{equation}\label{eq:incidence}
\begin{aligned}
\mathcal I=\{((Y,\alpha,[v])):&\ (\alpha X-B_0-YM)v=0\}\\
&\subseteq\mathbb C^{qa}\times\mathbb P^{1}\times\mathbb P^{d'-1}.
\end{aligned}
\end{equation}
Fix $(\alpha,[v])$ with $Mv\ne0$. The equation splits by rows as
$Y_i(Mv)=((\alpha X-B_0)_i)v$, one nontrivial linear equation on each $Y_i\in\mathbb C^{a}$, with
disjoint sets of unknowns; hence the fibre of $\mathcal I$ over $(\alpha,[v])$ has dimension
$qa-q$, and $\dim\mathcal I\le qa-q+1+(d'-1)=qa-q+d'<qa$ because $q>d'$. Therefore the projection of
$\mathcal I$ to $\mathbb C^{qa}$ is contained in a proper closed subset, and for $Y$ in a dense open
set no common root has $Mv\ne0$: every common root satisfies $Mv=0$, and then the equation reduces to
$(\alpha X-B_0)v=0$, which does not involve $Y$. Hence all common roots come from the finite
parameter-free set $\{\alpha:\operatorname{rank}(\alpha X-B_0)<d'\}$, and $g\in\mathbb C[\lambda]$.
\end{proof}

\begin{thm}[Closed and thin]\label{thm:closed}
$\mathcal B$ is closed in $\mathrm{Gr}(d',k)$; the feasible set is relatively open in the
compatibility stratum $\mathcal C=\{U:\ \mathcal F_U\ne\varnothing\}$, which is itself closed.
Moreover $\mathcal B$ has empty interior.
\end{thm}

\begin{proof}
Badness is the vanishing of the coefficient polynomials $q_{I,m}$ of Theorem~\ref{thm:rp} read as
functions of the Plücker coordinates: the gcd of a bounded-degree family of univariate polynomials is
one exactly when a certain Bezout system has a solution, which is an open condition on the
coefficients; its complement is closed. Closedness of $\mathcal C$ is bilinearity of
$UZ''\subseteq U+T$. Emptiness of the interior follows from the example of
Section~\ref{sec:examples}, where a bad point has feasible points arbitrarily close.
\end{proof}

\begin{thm}[Infeasibility]\label{thm:no}
Let $\mathcal C$ be the compatibility stratum for the given $d'$ and $\mathcal B$ the bad set. There
is no feasible $U$ of dimension $d'$ if and only if $\mathcal C\subseteq\mathcal B$. If
$\mathcal C\not\subseteq\mathcal B$ then the feasible set is dense in $\mathcal C$.
\end{thm}

\begin{proof}
Immediate from Theorem~\ref{thm:closed}: the feasible set is $\mathcal C\setminus\mathcal B$, and a
nonempty relatively open subset of an irreducible closed set is dense. (For a reducible stratum,
apply the argument to each component.)
\end{proof}

We record a sufficient condition that is checkable in linear algebra and that explains the smallest
bad examples.

\begin{prop}[Column drop]\label{prop:drop}
Let $K=HZ''+G'$ and suppose that for some $\lambda_0$ and some set $J$ of columns with
$|J|>k-d'$ the columns $K_{:,j}=\lambda_0H_{:,j}$ hold identically in the fibre parameters for
$j\in J$. Then $U$ is bad.
\end{prop}

\begin{proof}
At $\lambda_0$, $P(\lambda_0)_{:,j}=\lambda_0H_{:,j}-K_{:,j}=0$ for $j\in J$, so at most $k-|J|$
columns of $P(\lambda_0)$ are nonzero and $\operatorname{rank}P(\lambda_0)\le k-|J|<d'$, for every
point of the fibre. Hence $(5b)$ fails everywhere.
\end{proof}

The converse of Proposition~\ref{prop:drop} is false; column dependencies of a more general shape,
and the bad examples of Section~\ref{sec:examples}, show that the mechanisms are not exhausted by
columns vanishing.

\section{Algorithms and examples}\label{sec:algorithms}

\begin{thm}[Region with automatic compatibility]\label{thm:regime}
Let $s=\dim(U_0\cap T)$ and suppose $d'+t-k\ge s$ with $t=\dim T$. Then, for every $d'$,
infeasibility of the $d'$-stratum is decidable by a one-sided randomised procedure: sample a rational
$U\in\mathrm{Gr}(d',k)$ containing $U_0$, test the fibre with Theorem~\ref{thm:rp}, and accept when a
witness is found.
\end{thm}

\begin{proof}
Under the hypothesis a generic $U$ satisfies $\dim(U+T)=k$, so compatibility $UZ''\subseteq U+T$ is
automatic and $\mathcal C$ contains a dense open subset of $\mathrm{Gr}(d',k)$. By
Theorem~\ref{thm:closed} the feasible set is open and dense whenever nonempty, and by
Theorem~\ref{thm:witness} a feasible point is certified by a witness; a random rational $U$ therefore
lands in the feasible set with probability bounded below by $1-2d'^{2}/|S|$ for a grid $S$, and the
one-sided test of Theorem~\ref{thm:rp} accepts. If the stratum is entirely bad no witness exists and
the procedure never accepts.
\end{proof}

\begin{thm}[The stratum $t=1$]\label{thm:t1}
If $t=1$ and $U$ is feasible then $U$ is invariant under $Z''$; consequently, for $t=1$ the existence
question is decided by enumerating the $d'$-dimensional invariant subspaces of $Z''$ that contain
$T$ and $U_0$, and testing each with Theorem~\ref{thm:criterion}.
\end{thm}

\begin{proof}
Feasibility forces $a=\dim(U\cap T)\ge1$ (otherwise the fibre is a single point and the pencil is
square, hence has an eigenvalue, so $(5b)$ fails); with $t=1$ this gives $T\subseteq U$, hence
$U+T=U$ and compatibility reads $UZ''\subseteq U$, i.e.\ invariance. The enumeration is finite and
each test is the criterion.
\end{proof}

\subsection{Examples}\label{sec:examples}

\begin{example}[A feasible fibre at $t=2$]\label{ex:feasible}
Take $k=3$, $q=3$, $d'=2$, $T=\operatorname{span}\{e_1,e_2\}$,
\[
H=\begin{bmatrix}2&2\\-1&-1\\2&1\end{bmatrix},\quad
Z''=\begin{bmatrix}1&2\\1&1\end{bmatrix},\quad U=\operatorname{row}H=\mathbb C^{2}.
\]
The fibre is $\mathcal F_U=T^{\,3}$ (three free parameters), and the three nonzero $2$-minors of
$\lambda H-HZ''-G'$ have gcd one, so the fibre is feasible by Theorem~\ref{thm:criterion}. The
corresponding observer is realisable over $\mathbb Q$ with $n=4$, $p=1$, $r=3$, and its conditions
(\ref{eq:fivea})--(\ref{eq:fiveb}) hold on the original matrices.
\end{example}

\begin{example}[A bad point with feasible neighbours]\label{ex:bad}
With $k=3$, $q=3$, $d'=2$, $T=\operatorname{span}\{e_1,e_2\}$,
\[
Z''=\begin{bmatrix}-2&-2&0\\-2&1&-2\\1&0&-1\end{bmatrix},\qquad
U=\{y:y_3=y_2\},
\]
the fibre is one-dimensional in each row and $\gcd=\lambda+3$: the whole fibre is bad. The
neighbouring spaces $U_s=\{y:y_3=y_2+s\,y_1\}$, $s\ne0$, are feasible, so the bad set has empty
interior --- the statement of Theorem~\ref{thm:closed}.
\end{example}

\begin{example}[The thin region can be feasible]\label{ex:thin}
$k=3$, $t=2$, $d'=1$, $q=2$, $T=\operatorname{span}\{e_1,e_2\}$, $U=U_0=\operatorname{span}\{e_1\}$,
$Z''=0$, $H$ with both rows $e_1$. Compatibility is automatic and the fibre is $U^{2}$; with
$G'_i=(y_i,0,0)$ the pencil is $\operatorname{diag}(\lambda-y_1,\lambda-y_2)$ on the first two
columns, whose $1$-minors have gcd $1$ for $y_1\ne y_2$. Hence the stratum is feasible although
$d'+t-k=0<s=1$; sampling heuristics that only look at $d'+t-k$ versus $s$ are not sound.
\end{example}

\section{Open problems}\label{sec:open}

\begin{enumerate}[label=(\Alph*),leftmargin=2.2em,itemsep=2pt]
\item \emph{The thin stratum.} When $d'+t-k<s$ and compatibility is a genuine constraint,
$\mathcal C$ is a proper subvariety and the sampling argument of Theorem~\ref{thm:regime} does not
apply. Deciding $\mathcal C\subseteq\mathcal B$ there --- equivalently deciding the existence
question --- is open. Example~\ref{ex:thin} shows that feasibility is possible in that region, so the
question is not vacuous.
\item \emph{Bit complexity.} Theorem~\ref{thm:criterion} is polynomial for fixed compressed rank in
the arithmetic model; its bit cost, and that of the coefficient test behind
Theorem~\ref{thm:closed}, are not accounted for here. The bound of Theorem~\ref{thm:witness} is on
the witness, not on the running time.
\item \emph{Algebraic witnesses.} Theorem~\ref{thm:witness} assumes $U$ rational. For $t=1$ the
feasible $U$ are invariant subspaces of $Z''$ and may be algebraic, of degree growing with $n$; then
the witness may be algebraic as well and its height is not bounded here.
\item \emph{Hardness.} No reduction is given in either direction. Since the region
$d'+t-k\ge s$ is decidable by a randomised procedure, a hardness proof would have to encode its
instance into the thin stratum (A).
\end{enumerate}

\subsection{Reproducibility and evidence}\label{sec:repro}

Every theorem of this paper is accompanied by evidence in the artifact repository. The Lean~4
development formalises the core of Theorem~\ref{thm:criterion} --- the specialisation step of
Lemma~\ref{lem:specialize}, in the form used to prove Corollary~\ref{cor:every-bad} --- together with
the fibre identity of Theorem~\ref{thm:fibre} and the rank inequality used in the proof of
Theorem~\ref{thm:factor}. The examples of Section~\ref{sec:examples} are certified by exact rational
scripts: the minors are computed symbolically over $\mathbb Q$, the gcd is taken in the polynomial
ring, and the conditions (\ref{eq:fivea})--(\ref{eq:fiveb}) are re-verified on the original matrices
for the realisable instances. The classification experiments that suggested
Theorem~\ref{thm:factor} and the counterexample of Example~\ref{ex:thin} are recorded as such; they
are not used as evidence for any statement in this paper.

\begin{thebibliography}{99}
\bibitem{fer10} J. A. Fernando, H. Trinh, L. Jennings. Functional observability and the design of minimum order linear functional observers. \emph{IEEE Trans. Automat. Control} 55(5):1268--1273, 2010.
\bibitem{dar25} M. Darouach, J. A. Fernando. On functional observability and functional observer design. \emph{Automatica}, vol. 173, Art. no. 112115, 2025, doi: 10.1016/j.automatica.2025.112115.
\bibitem{fer26a} J. A. Fernando. A minimum-order functional observer beyond the Darouach and Luenberger constructions. arXiv:2609.33615, 2026 (preprint).
\bibitem{fer26b} J. A. Fernando. Complete characterization of minimum-order functional observers from Darouach to Luenberger. arXiv:2609.37992, 2026 (preprint).
\bibitem{paperI} Z. Liu. Deciding functional observer feasibility: compression, certificates, and parameterized algorithms. Companion paper (Paper I).
\bibitem{paperII} Z. Liu. Minimum order and minimum-order construction for functional observers. Companion paper (Paper II).
\bibitem{hj13} R. A. Horn, C. R. Johnson. \emph{Matrix Analysis}. 2nd ed. Cambridge Univ. Press, 2013.
\bibitem{vzgg13} J. von zur Gathen, J. Gerhard. \emph{Modern Computer Algebra}. 3rd ed. Cambridge Univ. Press, 2013.
\bibitem{lan02} S. Lang. \emph{Algebra}. Revised 3rd ed. Springer, 2002.
\bibitem{emi99} I. Z. Emiris, B. Mourrain. Matrices in elimination theory. \emph{J. Symbolic Comput.} 28(1--2):3--44, 1999, doi: 10.1006/jsco.1998.0266.
\end{thebibliography}

\end{document}
